% ---
% id: receiver/rem:bound-data-band
% kind: rem
% title: "The bound-data test and the contractivity test on the band"
% summary: >-
%   Where the contractivity test at ε₁ admits and the bound-data test at ε₂
%   refuses, the bound-data test decides: an exact implication between the two
%   conditions transports no admission between computed numbers. Admitting the
%   band would make a stated error term false on admitted values.
% prominence: standard
% status: accepted
% uses: [def:tolmap, def:gram, def:bounddata, thm:multiclick-bound, thm:port-grams, thm:exact-and-evaluated, rem:cdirection, rem:tolerancedomain, cor:unconditioned-family]
% mentions: [rem:ranktwosided, ass:passive, receiver/component-catalogue/ex:bb84]
% assumes: []
% refs: []
% provenance: newly-derived
% lean:
% ---

Items (iv) and (v) of Definition~\ref{def:tolmap} both widen an exact
condition by a receiver tolerance, and in exact arithmetic the condition of
item (v) implies that of item (iv). This remark says why the implication
nevertheless admits nothing, and what each widened test does and does not
guard against.

\begin{remark}\label{rem:bound-data-band}
\emph{Why item (iv) is widened at all.} Admissible receivers with
$\lmax(B)=1$ exactly are ordinary: at a receiver without a probe $B=A_\Sigma$
(Corollary~\ref{cor:unconditioned-family}), and the channel-only receiver of
Example~\ref{ex:bb84} has $A_\Sigma=\id_2$. There a computed $\lmax(B)$ one
rounding unit above $1$ makes $c$ negative on a receiver that is admissible,
and refusing it would refuse the simplest legitimate device.

\emph{The two items decide different objects.} A declaration can satisfy one
and fail the other: item~(v) admits a route-(ii) $\Gmat$ at $\varepsilon_1$,
item (iv) decides the derived $c$ at $\varepsilon_2$, and whenever
$\varepsilon_1>\varepsilon_2$, the ordinary declaration, the band
\eqref{eq:tolmapband} is non-empty, admitted by item~(v)'s test and refused by
item (iv)'s. On route (i) the band has no analogue, item~(i) and item (iv)
carrying the same $\varepsilon_2$; but item (iv) is a test there too and not a
formality, for the reason step~(2) below gives. Item (iv) wins on the band, in
four steps.

(1) \emph{In exact arithmetic item~(v)'s condition implies item (iv)'s, so
item (iv) is redundant, exactly.} $\Gmat^\dagger\Gmat\preceq\id$ gives
$A_\Sigma\preceq\Gmat_{\Inp}^\dagger\Gmat_{\Inp}\preceq\id$ by
Theorem~\ref{thm:port-grams}\,(i), hence $\lmax(A_\Sigma)\le1$ and $c\ge0$; and at
every setting $B^{(t)}\preceq\id$ by \eqref{eq:bbkhyp}, hence $c\ge0$ again.

(2) \emph{But item~(v) never establishes its exact condition.} It admits under
a widened test on a \emph{computed} number, and Theorem~\ref{thm:exact-and-evaluated}\,(ii)
bounds the distance from that number to the exact one by nothing at all. So
what step~(1) transports is an implication between exact quantities and
\emph{no implication whatever between the two computed quantities}, which are
separated by at least one matrix product
($A_k=\Gmat_{\Inp}^\dagger\Emd{\{k\}}\Gmat_{\Inp}$,
Definition~\ref{def:gram}) and, at a setting of a family a submodule
constructs from the Grams, by the further operations of that construction as
well. Reading
item~(v) as admitting anything under item (iv) is reading an exact implication
as a bound on an evaluation, which is the one substitution
Theorem~\ref{thm:exact-and-evaluated}\,(ii) exists to forbid. This step is also why item (iv)
is not redundant on route (i), where step~(1) holds just as well from
$\sum_kA_k\preceq\id$: redundancy of the \emph{exact} condition is not
admission of the \emph{computed} one. And $0\le c$ is tested where the rest of
\eqref{eq:bounddata} is not because the error term of
Remark~\ref{rem:cdirection}, on which step~(3) turns, is stated for a negative
computed $c$ and bounded through this test alone.

(3) \emph{The band is not empty of consequence, and the consequence is
quantitative and unsafe.} Remark~\ref{rem:cdirection} bounds the even-$n$
overshoot of the unclamped evaluation by
$(K-1)|c|^{\,n}\le(K-1)\varepsilon_2^{\,2}$, and that bound is stated in
$\varepsilon_2$ \emph{because item (iv) is what bounds $|c|$}. Admitting
\eqref{eq:tolmapband} would make $|c|$ as large as $\varepsilon_1$ and the
true bound $(K-1)\varepsilon_1^{\,2}$, larger by the factor
$(\varepsilon_1/\varepsilon_2)^2$, which at an ordinary pair of declared
numbers is many orders of magnitude. A stated error term of this contract
would then be false on values this contract admits, in the direction where
the evaluated $\Bnd(n)$ \emph{exceeds} the exact one, which is the direction
that reaches a key length. Remark~\ref{rem:tolerancedomain} says what a
declared domain buys, ``that every error term this contract states holds on
every value the contract admits'', and that is the property item~(v) winning
would destroy.

(4) \emph{Nothing is lost by item (iv) deciding.} The case the widening exists
for, a receiver with $\lmax(B)=1$ exactly whose computed $\lmax(B)$ sits a
rounding unit above $1$, is admitted by item (iv)'s own $-\varepsilon_2$; and item (iv) tests the
very quantity that enters \eqref{eq:beta}, where item~(v) tests one two
computations upstream. A declarer who wants the wider band declares a wider
$\varepsilon_2$, which is what a receiver tolerance is for and which carries
the stated error term with it, rather than leaving it behind.

\emph{Item (v) is one-sided as a test and two-sided in its consequences.}
$\Gmat^\dagger\Gmat\preceq\id$ is an upper bound on a spectrum, so the test
relaxes in one direction only: it admits
$\lmax(\Gmat^\dagger\Gmat)\in(1,1+\varepsilon_1]$, which the exact condition
excludes, and it refuses nothing that the exact condition, applied to the same
computed number, would admit. The \emph{errors} are two-sided in the sense of
Remark~\ref{rem:ranktwosided}, and deliberately so. Too tight refuses a
legitimate device: a receiver with a unitary $\Gmat$, such as the worked
instrument \eqref{eq:bb84G} of Example~\ref{ex:bb84}, has
$\lmax(\Gmat^\dagger\Gmat)=1$ \emph{exactly}, so a computed value one rounding unit above
$1$ would refuse it, which is the argument for widening item~(iv) verbatim.
Too loose admits a genuinely non-contractive $\Gmat$, which violates
Assumption~\ref{ass:passive}. Route (ii) carries no separate Gram
admissibility check, because Theorem~\ref{thm:port-grams}\,(i) \emph{derives}
$\sum_kA_k\preceq\Gmat_{\Inp}^\dagger\Gmat_{\Inp}\preceq\id$ from
contractivity, so the \emph{exact} condition item~(i) tests at route (i) is
redundant here rather than dropped.

\emph{What that derivation transports, and what it does not.} It is an
implication between \emph{exact} quantities. It therefore makes item~(i)'s
exact condition redundant on route (ii), and, through $B\preceq\id$ in
\eqref{eq:bbkhyp}, item~(iv)'s exact condition $c\ge0$ redundant too. It
transports \emph{no admission}: the computed $\lmax(A_\Sigma)$ and the computed
$c$ are separate evaluations, related to the exact ones by nothing this
contract asserts (Theorem~\ref{thm:exact-and-evaluated}\,(ii)), so an $\varepsilon_1$-admitted
$\Gmat$ does not thereby satisfy any test stated at $\varepsilon_2$. Item~(iv)
accordingly decides the bound data on route (ii) as on every other, and the
band \eqref{eq:tolmapband} that opens when $\varepsilon_1>\varepsilon_2$ is
refused. \emph{Each route therefore tests its own declared datum under its own
declared number}, $\varepsilon_1$ on the supplied $\Gmat$ at route (ii) and
$\varepsilon_2$ on the supplied Grams at route (i), because each is the
tolerance attached to \emph{that} datum's residual, and neither is thereby a
tolerance on the derived $c$. A declarer who wants one slack on
$\lmax(A_\Sigma)$ regardless of route declares
$\varepsilon_1=\varepsilon_2$, which also empties \eqref{eq:tolmapband}.
\end{remark}
